\documentclass[11pt]{article}
\usepackage[margin=1in]{geometry}
\usepackage{amsmath,amssymb,amsthm}
\usepackage[hidelinks]{hyperref}
\newtheorem{theorem}{Theorem}
\begin{document}
\setlength{\emergencystretch}{1em}
\title{The ordinary degree-seven binomial residue frontier}
\author{}\date{}\maketitle
Let $\mathcal L(X^j)=j!$ and
\[
 P_j^{a,d}(t)=\sum_{k=0}^j\binom jk(aj+1)_{dk}t^k.
\]
Strong Factorial for a nonzero phase $f$ means that
$\mathcal L(f^m),\mathcal L(f^{m+1})$ cannot both vanish for any
positive integer $m$.

\begin{theorem}[Exact remaining degree-seven tests]
\label{thm:residue-frontier}
Strong Factorial for every ordinary complex univariate binomial
of degree at most seven is equivalent to
\[
 \gcd_{\mathbb Q[t]}(P_{N-1}^{a,d},P_N^{a,d})=1
\]
for every terminal integer $N\ge6$ in the following three rows:
\[
 \begin{array}{c|c|l}
 (a,d)&\text{support}&\text{terminal restriction}\\ \hline
 (1,6)&\{1,7\}&\text{all }N\ge6\\
 (3,4)&\{3,7\}&N\equiv14,46,55\pmod{64}\\
 (5,2)&\{5,7\}&N\equiv12,14\pmod{16}.
 \end{array}
\]
For the last row, the gcd condition is equivalent to excluding an
irreducible quadratic common factor.
No assertion of coprimality on all these remaining tests is made.
\end{theorem}
\begin{proof}
Corollary \texttt{cor:degree-seven-frontier} of the
\href{ORDINARY_DEGREE_SIX_BINOMIAL_SFC.tex}{degree-six assembly}
reduces the degree-seven assertion to the four supports
$\{1,7\},\{3,7\},\{4,7\},\{5,7\}$ at every terminal $N\ge6$.

Corollary \texttt{cor:four-seven} of the
\href{PRIME_PLUS_ONE_EXPONENT_GAP_SFC.tex}{prime-plus-one exponent note}
proves the support $\{4,7\}$ at every order.
Corollary \texttt{cor:sixty-one-sixty-four} of the
\href{LOWER_THREE_FOURTH_POWER_RESIDUE_REDUCTION.tex}{fourth-power residue note}
proves support $\{3,7\}$ at every terminal outside its displayed
three residue classes.
Theorem \texttt{thm:residues} of the
\href{LOWER_FIVE_QUADRATIC_RESIDUE_REDUCTION.tex}{quadratic residue note}
proves support $\{5,7\}$ at every terminal outside its displayed
two classes.
Theorem \texttt{thm:global-quadratic-bound} of that same note
proves that at every index its rational gcd is either constant or
an irreducible quadratic of negative discriminant.
This proves the additional equivalent criterion in the last row.
Thus the three rows are sufficient to complete all four-support
tests in the existing exact classification.
Conversely every test in the rows belongs to the degree-seven
assertion, so that assertion implies each row.

The ordinary moment normalization is
\[
 \mathcal L((X^a(A+BX^d))^j)=A^j(aj)!P_j^{a,d}(B/A)
 \qquad(AB\ne0).
\]
A common complex root of two rational polynomials is equivalent to
a nonconstant rational gcd. Nonzero degenerate monomials have all
positive factorial moments nonzero.
The imported classification already includes all degrees at most six
and all supports through terminal order five; no finite-prefix
inference is used.
\end{proof}

This is a reduction of the remaining problem.
For support $\{5,7\}$, every nonconstant rational common factor
is now confined to an irreducible imaginary quadratic at every index.
The remaining progressions, all-degree binomial SFC, and unrestricted
FC/SFC remain open. The written reduction and its input proofs are
internally checked; no independent review, formal verification or
priority is asserted.
\end{document}
